\section{Localization of the Hessian and of the fluctuation covariance}\label{sec:localization}

\paragraph{Status.} The localization of this section and of Section~\ref{sec:interaction-local} is of radius type (Section~\ref{sec:interaction-local}, Summary). The polymer expansion of Part~X needs volume type, which Part~IX provides. The main line uses from this section Lemma~\ref{lem:hessian-gamma}, Definition~\ref{def:dirichlet-hessian}, Proposition~\ref{prop:dirichlet-hessian}, Lemma~\ref{lem:coercive-general} and Proposition~\ref{prop:local-fluct}. Theorem~\ref{thm:local-covariance} and Corollary~\ref{cor:large-changes} (Section~\ref{subsec:local-cov-expansion}) are superseded by the weakened representations of Section~\ref{sec:nested-walks} and are not used later. Lemma~\ref{lem:telescoping} serves only these results and the radius-type results of Sections~\ref{sec:det-local} and~\ref{sec:interaction-local}, so it is not on the main line either.

\subsection{The telescoping principle}\label{ss:37.1}

For a finite set $\mathsf S$ of unit points and $m\ge0$ let $\Lam_m(\mathsf S)$ be the union of the blocks $B(y)$ with $\operatorname{dist}(y,\mathsf S)\le m$, and let $m_{\mathsf S}$ be the least $m$ with $\Lam_m(\mathsf S)=\Tf$.

\begin{lemma}[Telescoping]\label{lem:telescoping}
Let $F^X$ be numbers attached to the domains $X=\Lam_m(\mathsf S)$, $0\le m\le m_{\mathsf S}$, with $F^{\Tf}=F$ and $|F^{\Lam_m(\mathsf S)}-F^{\Lam_{m-1}(\mathsf S)}|\le Ke^{-\gamma m}$. Then
\[
F=\sum_{m=0}^{m_{\mathsf S}}F_m,\qquad F_0=F^{\Lam_0(\mathsf S)},\quad F_m=F^{\Lam_m(\mathsf S)}-F^{\Lam_{m-1}(\mathsf S)},\qquad|F_m|\le Ke^{-\gamma m}\ (m\ge1).
\]
If each $F^X$ is a function of the background that depends only on its restriction to a set $\mathcal R(X)$, is analytic, and is invariant (or covariant) under real gauge transformations, and if $\mathcal R(\Lam_{m-1}(\mathsf S))\subset\mathcal R(\Lam_m(\mathsf S))$ for $1\le m\le m_{\mathsf S}$, then $F_m$ has the same properties with $\mathcal R(\Lam_m(\mathsf S))$.
\end{lemma}

The proof is the identity of the telescoping sum. By the inclusion, $F^{\Lam_m(\mathsf S)}$ and $F^{\Lam_{m-1}(\mathsf S)}$ both depend only on the background on $\mathcal R(\Lam_m(\mathsf S))$, and so does $F_m$. Theorem~\ref{thm:local-covariance}, Proposition~\ref{prop:det-local} and Theorem~\ref{thm:interaction-local} are applications of this lemma; Corollary~\ref{cor:composed} uses the same telescoping identity directly. In these three applications $\mathcal R(X)$ is $\widehat{X^+}$ (Theorem~\ref{thm:local-covariance}, Proposition~\ref{prop:det-local}) or $\hat X$ together with the bonds of $\mathfrak B(X)$ for the coarse field (Theorem~\ref{thm:interaction-local}); these sets increase with $X$, so the inclusion holds. The estimates needed are comparison bounds for approximants on nested domains, whose differences have seam form.

\subsection{The coarse Hessian as an operator of Part~II}\label{ss:37.2}

In this section $W$ has two meanings, each fixed where it occurs: the base configuration $W=U^{(k)}$ of the step (Section~\ref{ss:31.3}), and a potential, as in $\bG[W]$, $W_2$ and $W_2^X$.

Let $\bH_0=\bH[0]$ be the minimizer map without potential, and for a regular background $U$ put
\[
\omega_0=\bH_0^TJ_U,\qquad\ip A{W_2A}=d^2\mathcal N(0)[A,A]-\ip{\omega_0}{d^2\Cav_k(0)[A,A]},
\]
with $\mathcal N$ the nonlinear part of the action at $U$ (Proposition~\ref{prop:action}) and $\Cav_k$ that of the gauge-adapted averaging (Proposition~\ref{prop:Q-expansion}).

\begin{lemma}[{The Hessian is $\boldsymbol\Gamma_\pi[W_2]-a$}]\label{lem:hessian-gamma}
$W_2$ is a real symmetric bounded operator with $\|W_2\|\le C\eps$ and $\|\one_yW_2\one_{y'}\|\le C\eps e^{-\gamma|y-y'|}$. If $U$ is critical, i.e.\ $U=U'$ is the background of the $k$-step problem for $W=U^{(k)}$, then $\omega_0$ is its multiplier and
\[
\Delta^0=\boldsymbol\Gamma_\pi[W_2]-a,
\]
where $\Delta^0$ is the Hessian of Proposition~\ref{prop:hessian} and $\boldsymbol\Gamma_\pi[W_2]$ is the operator of Theorem~\ref{thm:reg-min} with the potential $W_2$.
\end{lemma}

\begin{proof}\hasdetail{lem:hessian-gamma}
At $A=0$ the second derivative $d^2\mathcal N(0)$ is a sum of plaquette-local bilinear forms. By the first part of Lemma~\ref{lem:div-curvature}, it is a local operator with $|d^2\mathcal N(0)[h,h']|\le C\eps\|h\|\,\|h'\|$; its proof is the computation of the proof of Proposition~\ref{prop:action}(c) at $A=0$, with the lattice trick \eqref{eq:lattice-trick}. By Proposition~\ref{prop:Q-expansion}(c), $|d^2\Cav_k(0)[h,h](c)|\le C\|h\|^2_{S_c}$. By Cauchy--Schwarz on blocks, as in the proof of Proposition~\ref{prop:dirichlet-hessian}(b), $|\omega_0(c)|\le\sum_{y'}\|\one_{y'}\bH_0\one_c\|\,\|J_U\|_{B(y')}\le C\eps$, by the block decay of $\bH_0$ (Theorem~\ref{thm:reg-min}) and $|J_U|\le\eps$. Hence the bounds. The coefficients $\omega_0(c)$ depend on $U$ through the decaying kernel of $\bH_0^T$.

If $U$ is critical, the Euler--Lagrange equation at $B=0$ gives $J=\bQ^T\omega$. So $\bH_0^TJ=(\bQ\bH_0)^T\omega=\omega$. By Proposition~\ref{prop:hessian}(a), $\Delta^0(B,B)=\min\{\|DA\|^2+\ip A{W_2A}:\bQ A=B,\ RD^TA=0\}$. By Theorem~\ref{thm:cov-min}(b) with the potential $W_2$, the minimizer is $\bH[W_2]B$. On the constraint set $\ip A{\fA_\pi A}=\|DA\|^2+\ip A{W_2A}+a\|B\|^2$. The minimum of $\ip A{\fA_\pi A}$ over $\{\bQ A=B\}$ is $\ip B{\boldsymbol\Gamma_\pi B}$ and is attained at $\bH[W_2]B$. Only the $L^2$ properties of $W_2$ enter: the coercivity of Theorem~\ref{thm:reg-coercive} uses $W$ only through $\|W\|$, and the variational characterization is algebraic.
\end{proof}

\subsection{Dirichlet Hessians}\label{ss:37.3}

\begin{definition}[Dirichlet Hessian]\label{def:dirichlet-hessian}
For a domain $X$ put
\[
\begin{aligned}
\omega_0^X&=(\bH_0^X)^T\one_XJ_U\quad\text{on }\mathfrak B(X),\\
\ip A{W_2^XA}&=d^2\mathcal N(0)[\one_XA,\one_XA]-\ip{\omega_0^X}{d^2\Cav_k(0)[A,A]}_{\mathfrak B(X)},
\end{aligned}
\]
for $A\in\ell^2_1(X)$, and
\[
\Delta^{0,X}=\boldsymbol\Gamma^X[W_2^X]-a\qquad\text{on }\ell^2(\mathfrak B(X)).
\]
\end{definition}

\begin{proposition}[Dirichlet Hessians]\label{prop:dirichlet-hessian}
\begin{enumerate}[label=(\alph*)]
\item For $X=\Tf$: $\omega_0^{\Tf}=\omega_0$ and $\Delta^{0,\Tf}=\boldsymbol\Gamma_\pi[W_2]-a$. This is $\Delta^0$ if $U$ is critical.
\item $|\omega_0^X|\le C\eps$, $\|W_2^X\|\le C\eps$, and $\Delta^{0,X}$ has decay $(C,\gamma)$, uniformly in $X$.
\item $\omega_0^X$, $W_2^X$ and $\Delta^{0,X}$ depend only on $U$ on $\hat X$ and are covariant under real gauge transformations. They are analytic under complex perturbations of $U$ as in Proposition~\ref{prop:reg-complex}. The potential $W_2^X$ is real symmetric at real $U$ and has range $\le3$ in $A$, but it depends on $U$ on all of $\hat X$. It satisfies the hypotheses of Proposition~\ref{prop:reg-complex} for existence, $\Rea$-coercivity and decay; its comparison of domains is (d). The analyticity and the changes under complex perturbations come from:
  \begin{itemize}
  \item $d^2\mathcal N(0)$ and the kernel of $d^2\Cav_k(0)$ are local and analytic in the background (Propositions~\ref{prop:action} and~\ref{prop:Q-expansion}, and Proposition~\ref{prop:AB} at complex references);
  \item $\ip b{\omega^X_0}=\ip{F^s_U}{D\bH^X_0b}$, with $D\bH^X_0$ analytic by Proposition~\ref{prop:reg-complex} at $W=0$;
  \item the changes of size $C\delta$ follow from Cauchy's estimate on a ball of twice the radius.
  \end{itemize}
\item For $X'\subset X$ and $c,c'\in\mathfrak B(X')$,
\[
\begin{aligned}
|\omega_0^X(c)-\omega_0^{X'}(c)|&\le C\eps\sum_{w\in\Sigma}e^{-\gamma|c-w|},\\
|\Delta^{0,X}(c,c')-\Delta^{0,X'}(c,c')|&\le C\sum_{w\in\Sigma}e^{-\gamma'(|c-w|+|w-c'|)},
\end{aligned}
\]
with $\Sigma$ as in Proposition~\ref{prop:dirichlet-averaged}(d) and some $\gamma'>0$.
\end{enumerate}
\end{proposition}

\begin{proof}\hasdetail{prop:dirichlet-hessian}
(a)--(c) follow from Proposition~\ref{prop:dirichlet-averaged} and Lemma~\ref{lem:hessian-gamma}. The bound $|\omega_0^X|\le C\eps$ needs no pointwise estimate. By Cauchy--Schwarz on blocks, $|\omega_0^X(c)|\le\sum_{y'}\|\one_{y'}\bH_0^X\one_c\|\,\|J_U\|_{B(y')}\le C\eps\sum_{y'}e^{-\gamma|c-y'|}$, using only the block kernel bounds of $\bH^X_0$ and $|J_U|\le\eps$. (d) The bound for $\omega_0$ follows from Proposition~\ref{prop:dirichlet-averaged}(d) for $\bH_0$, since $J_U$ is local and $|J_U|\le\eps$. Then $W_2^X-W_2^{X'}$ on $X'$ is bounded near $w$ by $C\eps\sum_{s\in\Sigma}e^{-\gamma|w-s|}$. The $d^2\mathcal N$ parts differ only on bonds within distance $1$ of $X\setminus X'$. Insert this into Proposition~\ref{prop:dirichlet-averaged}(d) for $\boldsymbol\Gamma$. Since $\sum_we^{-\gamma(|c-w|+|w-c'|)}e^{-\gamma|w-s|}\le Ce^{-\gamma'(|c-s|+|s-c'|)}$, the result has seam form.

\emph{Complex perturbations in (c).} Let $U_\delta$ be as in Proposition~\ref{prop:reg-complex}, of size $\delta<2\delta_*$, with $\delta_*$ fixed at the end of this paragraph. In a local real gauge of $U$, $U_\delta=e^{\ii\eta A_0}$ with $A_0$ satisfying (B) with parameter $C(\eps+\delta)$ (Section~\ref{sec:complex-references}). This uses the bound on $\partial\delta\!A$: with a sup bound alone, $H_p-1$ would only be $O(\eta\delta)$. Hence $|H_p(U_\delta)-1|\le C\eta^2(\eps+\delta)$, $\|R(U_\delta(b))^{\pm1}\|\le2$ and $|F^s_{U_\delta}|\le C(\eps+\delta)$ (step~(v) of the proof of Proposition~\ref{prop:complex-ext}). The proof of the bound on $W_0$ in Lemma~\ref{lem:div-curvature} uses only these three bounds, the transports being bounded by $2$ instead of $1$. So $d^2\mathcal N_{U_\delta}(0)$ is a local operator of $L^2$ norm at most $C(\eps+\delta)$. Read $\omega_0^X$ through the identity $\ip b{\omega_0^X}=\ip{F^s}{D\bH^X_0b}$ of (c), and use Proposition~\ref{prop:AB}(iii),(vi) for $d^2\Cav_k(0)$. Then $W_2^X(U_\delta)$ is holomorphic in $\delta\!A$, of range $\le3$, with $\|W_2^X(U_\delta)\|\le C(\eps+\delta)$, for every $\delta\!A$ of size $<2\delta_*$. The Schwarz lemma on the complex lines $t\mapsto U_{t\,\delta\!A}$, $|t|<2\delta_*/\delta$, gives $\|W_2^X(U_\delta)-W_2^X(U)\|\le C(\eps+2\delta_*)\,\delta/\delta_*$. For $\delta\le\delta_*$ this is at most $C(\eps+2\delta_*)$. So $W_2^X$ satisfies the range-$3$ clause of Proposition~\ref{prop:reg-complex} at these perturbations once $C'(\eps+2\delta_*)\le c$, with $c$ the stability threshold of Lemma~\ref{lem:CT-stable}. Here $\delta_*>0$ is fixed once, depending on $d$, $L$, $G$, $\alpha$ and $M$, for all $\eps$ with $\eps\le\frac12\min(\eps_0,c'_3)$, $C\eps\le\frac12c_5$ and $C'\eps\le\frac12c$. It satisfies $2C\delta_*\le\frac12\min(\eps_0,c'_3)$, $2C\delta_*\le\frac12c_5$, $2C'\delta_*\le\frac12c$, and the radius conditions of Propositions~\ref{prop:reg-complex} (with $W=0$), \ref{prop:AB} and~\ref{prop:dirichlet-averaged}(e) at size $2\delta_*$. With $C'\eps\le\frac12c$, a threshold of this section, the clause holds for all $\delta\le\delta_*$. Moreover $\|W_2^X(U_\delta)-W_2^X(U)\|\le C_W\delta$ with $C_W=C(1+2\delta_*)/\delta_*$, a fixed constant. The constant $c''_L$ of Lemma~\ref{lem:coercive-general} is chosen with $C_{\rm cmp}c''_L\le\delta_*$, with $C_{\rm cmp}=4$ of Section~\ref{sec:complex-references}.
\end{proof}

\subsection{The fluctuation operator on domains}\label{ss:37.4}

We return to the step of Part~VII. Its operator $\Delta_k$ acts on fields on the set $\mathfrak F$ of bonds of $\T$ other than the bonds $b_0(c)$ (Definition~\ref{def:bonds}). By Lemma~\ref{lem:step-quadratic} it consists of two parts:
\begin{itemize}
\item the coarse-action part $\Delta^0(\hat Z^1\cdot,\hat Z^1\cdot)+2\ip{\omega_0}{\hat Z^2}$, built with the elimination maps $\hat Z^1$, $\hat Z^2$ of Lemma~\ref{lem:eliminate}. These depend only on $W=U^{(k)}$ within distance $2L$, and $\hat Z^1\delta_b$ is supported within distance $2L$ of $b$;
\item the gauge-fixing part $\ip{\mathfrak G_WZ}{\mathfrak G_WZ}$, with $\mathfrak G_W$ the linearized averaged contour variables at the base (Definition~\ref{def:avg-transport}). By Lemma~\ref{lem:avg-transport}(a), $\mathfrak G_WZ(y,x)$ depends only on $Z$ and $W$ on the inner bonds of the block $B'(y)$. So the kernel of $\mathfrak G_W^T\mathfrak G_W$ on fields on $\mathfrak F$ vanishes unless $b$ and $b'$ are inner bonds of a common block; it is block-diagonal, with range at most $L$.
\end{itemize}
For a domain $X$ let $\T'(X)$ be the set of the $y\in\T'$ with $B'(y)\subset X$, and put
\[
\ip Z{K^XZ'}=\sum_{y\in\T'(X)}\ \sum_{x\in B'(y)}\ip{\mathfrak G_WZ(y,x)}{\mathfrak G_WZ'(y,x)},\qquad\text{i.e.}\quad K^X=\mathfrak G_W^T\one_{\T'(X)}\mathfrak G_W ,
\]
on fields on $\mathfrak F$, with the bilinear pairing of Section~\ref{sec:fluctuations}. A block contributes to $K^X$ if and only if it lies in $X$; blocks that meet $\Tf\setminus X$ contribute nothing. For $X=\Tf$, $K^{\Tf}=\mathfrak G_W^T\mathfrak G_W$. For a domain $X$ and $b,b'\in\mathfrak F$ define the \emph{Dirichlet entries}
\[
\Delta^{(X)}_k(b,b')=\Delta^{\rm c,X}(b,b')+K^X(b,b'),
\]
where $\Delta^{\rm c,X}(b,b')=\Delta^{0,X}\big(\hat Z^1\delta_b,\hat Z^1\delta_{b'}\big)+\big(\text{polarization of }Z\mapsto2\ip{\omega_0^X}{\hat Z^2(Z)}\big)(b,b')$ if $\hat Z^1\delta_b$ and $\hat Z^1\delta_{b'}$ are supported in $\mathfrak B(X)$, and $\Delta^{\rm c,X}(b,b')=0$ otherwise. For $X=\Tf$ these are the entries of $\Delta_k$ (Lemma~\ref{lem:step-quadratic}), at critical $U$ (Proposition~\ref{prop:dirichlet-hessian}(a)). The entries $K^X(b,b')$ depend only on $U$ on $\hat X$, are covariant under real gauge transformations (Lemma~\ref{lem:avg-transport}(d)), and are analytic under the complex perturbations of $U$ considered below (Lemma~\ref{lem:avg-transport}(a),(f)).

Fix a margin $R_0\ge2L+2$. For a domain $Y$ let $Y^+=Y^{(R_0)}$, let $\mathfrak F(Y)$ be the set of bonds of $\mathfrak F$ starting in $Y$, and define the \emph{local fluctuation operator}
\[
\Delta_Y=\one_{\mathfrak F(Y)}\Delta^{(Y^+)}_k\one_{\mathfrak F(Y)}\qquad\text{on }\ell^2(\mathfrak F(Y)).
\]
An inner bond starting in $Y$ lies in a block $B'(y)\subset Y^{(L)}\subset Y^+$. Hence the gauge-fixing part of $\Delta_Y$ is the compression of that of the torus,
\begin{equation}\label{eq:gf-local-part}
\one_{\mathfrak F(Y)}K^{Y^+}\one_{\mathfrak F(Y)}=\one_{\mathfrak F(Y)}\mathfrak G_W^T\mathfrak G_W\one_{\mathfrak F(Y)} ,
\end{equation}
and only the coarse-action part of $\Delta_Y$ depends on the Dirichlet construction.

For a regular background $U$ that is not necessarily critical, we \emph{define} $\Delta_k=\Delta^{(\Tf)}_k$ and $C_k=\Delta_k^{-1}$, with the elimination at the base $W=U^{(k)}$. This agrees with Part~VII when $U$ is critical (Proposition~\ref{prop:dirichlet-hessian}(a)).

\begin{lemma}[Coercivity at general backgrounds]\label{lem:coercive-general}
Under the smallness hypotheses of Corollary~\ref{cor:covariance}, but for any real regular $U$ with $|J_U|\le\eps$ (critical or not), $\frac14c_L\le\Delta_k\le C_L$ on $\ell^2(\mathfrak F)$. There is $c''_L>0$, depending on $L$ and $M$, with $C_{\rm AB}c''_L\le c'_L$ for the threshold $c'_L$ of Corollary~\ref{cor:covariance}(c), such that $\Rea\Delta_k\ge\frac18c_L$ for complex perturbations of size at most $c''_L$ in the local gauges, if moreover $\eps\le\eps_{\rm cc}$ and $\eps$ is below the thresholds of the proof of Proposition~\ref{prop:dirichlet-hessian}(c).
\end{lemma}

\begin{proof}\hasdetail{lem:coercive-general}
\emph{Real backgrounds.} Let $Z$ be a field on $\mathfrak F$ and $B=\hat Z^1Z$. Then $B\in\mathcal Z=\ker l^{\rm pl}$, $\|B\|\ge\|Z\|$, and $\|\mathfrak G_WZ\|=\|\mathfrak G_WB\|$ (Lemmas~\ref{lem:eliminate}(d) and~\ref{lem:step-quadratic}). Theorem~\ref{thm:step-coercive} and Lemma~\ref{lem:local-comparison} do not use criticality. Their proofs use only the $(M,\eps)$-regularity of the reference, through Lemma~\ref{lem:large-cubes}, which needs $\eps\le\eps_{\rm cmp}(M'_L)$ (part of the smallness hypotheses of Corollary~\ref{cor:covariance}), and Lemma~\ref{lem:avg-local-gauge}, and the expansions of Part~IV at the base $W=U^{(k)}$, which hold at every real base of small curvature (Section~\ref{sec:adapted}). So they give
\[
\ip B{(\boldsymbol\Gamma_\pi[0]-a)B}+\|\mathfrak G_WB\|^2\ge c_L\|B\|^2\qquad(B\in\mathcal Z).
\]
By Lemma~\ref{lem:hessian-gamma}, $\|W_2\|\le C\eps$. Lemma~\ref{lem:perturbed-inverse}, applied to $\bQ\bG[W]\bQ^T$, gives $\|\boldsymbol\Gamma_\pi[W_2]-\boldsymbol\Gamma_\pi[0]\|\le C\eps$. Since $\Delta^{0,\Tf}=\boldsymbol\Gamma_\pi[W_2]-a$ (Proposition~\ref{prop:dirichlet-hessian}(a)),
\[
\Delta^{0,\Tf}(B,B)+\|\mathfrak G_WB\|^2\ge(c_L-C\eps)\|B\|^2\ge\tfrac12c_L\|Z\|^2
\]
for $\eps$ small. The term $2\ip{\omega_0}{\hat Z^2(Z)}$ is bounded by $CL^{3d+1}\eps\|Z\|^2$, since $|\omega_0|\le C\eps$ (Lemma~\ref{lem:eliminate}(d)); this is at most $\frac14c_L\|Z\|^2$ by the threshold $CL^{3d+1}\eps\le\frac14c_L$ of Corollary~\ref{cor:covariance}. This gives the lower bound. The upper bound follows from Proposition~\ref{prop:dirichlet-hessian}(b), the locality of $\hat Z^1$, $\hat Z^2$, and $\|\mathfrak G_W\|\le C_L$ (Lemma~\ref{lem:avg-transport}(c)).

\emph{Complex perturbations.} Let $U=e^{\ii\eta A'}U_r$ be a complex reference of size at most $c''_L$ (Section~\ref{sec:complex-references}). Since $C_{\rm cmp}c''_L\le\delta_*$, it is a perturbation $U_\delta$ of $U_r$ as in Proposition~\ref{prop:reg-complex} with $\delta\le\delta_*$. The changes of the terms of $\Delta_k(Z,\bar Z)$ are bounded as follows.
\begin{itemize}
\item The operators of Parts~II, III and~VIII change by $O(c''_L)$ in norm (Proposition~\ref{prop:reg-complex}, and for $W_2$ the proof of Proposition~\ref{prop:dirichlet-hessian}(c)); so $\Delta^{0,\Tf}$ changes by $O(c''_L)$.
\item $|\omega_0|\le C(\eps+c''_L)$ (Proposition~\ref{prop:hessian-general}(b)), so the term with $\omega_0$ changes by at most $CL^{3d+1}c''_L\|Z\|^2$.
\item The base $W_k(U)$ is $e^{\ii\beta}W$ bondwise, with $W=W_k(U_r)$ and $|\beta|_\infty\le C_{\rm AB}c''_L$ (Proposition~\ref{prop:AB}(i)). Since $C_{\rm AB}c''_L\le c'_L$, the threshold of Corollary~\ref{cor:covariance}(c), the elimination at this base is analytic in $\beta$, with the bounds of Lemma~\ref{lem:eliminate}, by the argument in the proof of that corollary. So $\hat Z^1$ and $\hat Z^2$ change by local maps of size $C_Lc''_L$.
\item The gauge-fixing term. $W$ has curvature at most $2C_{\rm curv}\eps$ in the units of step $k$ (proof of Corollary~\ref{cor:parts-apply}, and Proposition~\ref{prop:avg-regular}). For $\eps$ and $c''_L$ small depending on $L$, Lemma~\ref{lem:avg-transport}(f) applies in the units of step $k$, i.e.\ with $\delta=1$. It gives $|(\mathfrak G_{W_k(U)}-\mathfrak G_W)Z(y,x)|\le CL^2|\beta|_\infty\max|Z|$, the maximum over the bonds of $B'(y)$. Summing over the $L^d$ points $x\in B'(y)$, the block-local operator $\mathfrak G_{W_k(U)}-\mathfrak G_W$ has norm at most $CL^{d/2+2}|\beta|_\infty$, while $\|\mathfrak G_W\|\le CL^{d/2+1}$ by Lemma~\ref{lem:avg-transport}(c). So $\ip{\mathfrak G_{W_k(U)}Z}{\mathfrak G_{W_k(U)}\bar Z}$ differs from $\|\mathfrak G_WZ\|^2$ by at most $CL^{d+3}C_{\rm AB}c''_L\|Z\|^2$.
\end{itemize}
Hence $\Rea\Delta_k(Z,\bar Z)\ge\big(\frac14c_L-C''_Lc''_L\big)\|Z\|^2$, with a constant $C''_L\ge C_L$ absorbing $C_W$ of the proof of Proposition~\ref{prop:dirichlet-hessian}(c). The largest explicit power of $L$ comes from the term with $\omega_0$. So $c''_L$ is chosen with $C_{\rm AB}c''_L\le c'_L$, with $C_{\rm cmp}c''_L\le\delta_*$ (proof of Proposition~\ref{prop:dirichlet-hessian}(c)), and so small that the sum $C''_Lc''_L$ of the four changes is at most $\frac18c_L$; the term with $\omega_0$ has the largest power $L^{3d+1}$, in $d=4$ a factor $L^{-13}$. Then $\Rea\Delta_k\ge\frac18c_L$.
\end{proof}

\begin{proposition}[Local fluctuation operators]\label{prop:local-fluct}
There is $R_0(L,M)\le C_LM$ such that for $R_0\ge R_0(L,M)$, under the hypotheses of Lemma~\ref{lem:coercive-general}:
\begin{enumerate}[label=(\alph*)]
\item $\Delta_{\Tf}=\Delta_k$, and $\|\Delta_Y-\one_{\mathfrak F(Y)}\Delta_k\one_{\mathfrak F(Y)}\|\le C_Le^{-\gamma'(R_0-2L-10M)/2}$, with $\gamma'$ of Proposition~\ref{prop:dirichlet-hessian}(d), taken valid also at the complex perturbations of Lemma~\ref{lem:coercive-general}. Hence $\Delta_Y\ge\frac18c_L$, and $C_Y=\Delta_Y^{-1}$ has decay $(C_L,\gamma_C)$ uniformly in $Y$.
\item $\Delta_Y$ and $C_Y$ depend only on $U$ on $\widehat{Y^+}$. They are covariant, and analytic under the complex perturbations of Lemma~\ref{lem:coercive-general}, with $\Rea\Delta_Y\ge\frac1{16}c_L$. At complex perturbations the current $J_U$ is not bounded uniformly in $\eta$. Wherever $\omega_0$ enters, write $\ip b{\omega_0^X}=\ip{F^s_U}{D\bH_0^Xb}$ and use the decay of $D\bH_0^X$, as in Section~\ref{sec:nested-walks}, level~(L4); the proofs of Proposition~\ref{prop:dirichlet-hessian}(b),(c),(d) and Lemma~\ref{lem:coercive-general} are to be read this way. In the proof of Proposition~\ref{prop:dirichlet-hessian}(d) this gives $\ip b{\omega_0^X-\omega_0^{X'}}=\ip{F^s_U}{D(\bH_0^X-\bH_0^{X'})b}$, which is of seam form, at all fine unit points, by Proposition~\ref{prop:dirichlet-averaged}(d),(e) for $D\bH_0$ at the complex perturbation.
\item For $Y\subset Y'$ and $b,b'\in\mathfrak F(Y)$,
\[
|(C_{Y'}-C_Y)(b,b')|+|(\Delta_{Y'}-\Delta_Y)(b,b')|\le C_L\sum_{w\in\Sigma_Y}e^{-\gamma''(|b-w|+|w-b'|)},
\]
where $\Sigma_Y$ is the set of unit points within distance $R_0+10M+2L$ of $Y'\setminus Y$.
\item \emph{(Uniformity in $r\ge0$.)} $\Rea(\Delta_Y+r)\ge\frac1{16}c_L+r$, and $(\Delta_Y+r)^{-1}$ has decay $\big(C_L(\frac1{16}c_L+r)^{-1},\gamma_C\big)$ uniformly in $r$ and $Y$. Also (c) holds for $(\Delta_Y+r)^{-1}$, with an extra factor $(\frac1{16}c_L+r)^{-1}$ on the right. Indeed, $r$ cancels from the operator of Lemma~\ref{lem:compare-inverse}, and conjugation does not see it, since $(K+r)_\alpha-(K+r)=K_\alpha-K$. This is the form used at level~(L5) of Section~\ref{sec:nested-walks} and in Proposition~\ref{prop:det-local}.
\end{enumerate}
\end{proposition}

\begin{proof}\hasdetail{prop:local-fluct}
(a) For $b,b'\in\mathfrak F(Y)$ the supports of $\hat Z^1\delta_b$, $\hat Z^1\delta_{b'}$ lie at distance at least $R_0-2L$ from $\Tf\setminus Y^+$. The gauge-fixing parts of $\Delta_k(b,b')$ and $\Delta^{(Y^+)}_k(b,b')$ agree, by \eqref{eq:gf-local-part}. By Proposition~\ref{prop:dirichlet-hessian}(d) with $X=\Tf$ and $X'=Y^+$, together with the locality of $\hat Z^1$, $\hat Z^2$,
\[
|\Delta_k(b,b')-\Delta^{(Y^+)}_k(b,b')|\le C_Le^{-\gamma'(R_0-2L-10M)/2}e^{-\gamma'|b-b'|/2}.
\]
Here the sum over the points $w$ of the seam is taken with half of the rate in the distance to the seam, which absorbs the number of these points. A Schur bound gives the norm estimate, with a constant $C_{\rm seam}$. By the same proof, read as in (b), the bound holds also at the complex perturbations of Lemma~\ref{lem:coercive-general}; we take $C_{\rm seam}$ and $\gamma'$ valid at real $U$ and at these perturbations. Put $R_0(L,M):=2L+10M+\frac2{\gamma'}\log^+\frac{16C_{\rm seam}}{c_L}$. For $R_0\ge R_0(L,M)$ the norm in (a) is at most $\frac1{16}c_L$, at real $U$ and at these complex perturbations. The constants $C_{\rm seam}$, $c_L$ and $\gamma'$ do not depend on $M$, so $R_0(L,M)\le C_LM$ for $M\ge1$; and $R_0(L,M)\ge2L+2$. Then $\one_{\mathfrak F(Y)}\Delta_k\one_{\mathfrak F(Y)}\ge\frac14c_L$, as a compression of $\Delta_k$ (Lemma~\ref{lem:coercive-general}). This gives the lower bound, since $\frac14-\frac1{16}\ge\frac18$, and Lemma~\ref{lem:CT-bounded} gives the decay.

(b) follows from Proposition~\ref{prop:dirichlet-hessian}(c), the locality of the elimination, the properties of $K^X$ stated before the proposition, and (a) for complex perturbations: there $\Rea\one_{\mathfrak F(Y)}\Delta_k\one_{\mathfrak F(Y)}\ge\frac18c_L$ by Lemma~\ref{lem:coercive-general}, so $\Rea\Delta_Y\ge\frac18c_L-\frac1{16}c_L$.

(c) Apply Lemma~\ref{lem:compare-inverse} with $\fA=\Delta_{Y'}$ and $\fA'=\Delta_Y$. The operator $\Delta=\Delta_{Y'}\one_{\mathfrak F(Y)}-\one_{\mathfrak F(Y)}\Delta_Y\one_{\mathfrak F(Y)}$ has two kinds of entries. For $b\in\mathfrak F(Y)$ they are $\Delta^{(Y'^+)}_k(b,b')-\Delta^{(Y^+)}_k(b,b')$. Their gauge-fixing parts cancel, since both equal $(\mathfrak G_W^T\mathfrak G_W)(b,b')$: the block of an inner bond starting in $Y$ lies in $Y^+\subset Y'^+$. The remaining coarse-action parts are of seam form by Proposition~\ref{prop:dirichlet-hessian}(d), since $Y'^+\setminus Y^+$ lies within distance $R_0$ of $Y'\setminus Y$. For $b\in\mathfrak F(Y')\setminus\mathfrak F(Y)$ they are entries $\Delta^{(Y'^+)}_k(b,b')$ with $b$ starting in $Y'\setminus Y$, bounded by $C_Le^{-\gamma_1|b-b'|}$, since the gauge-fixing part has range at most $L$. Both have seam form, and the decay of $C_{Y'}$ and $C_Y$ gives the bound for $C$. The bound for $\Delta$ is the first kind of entry.
\end{proof}

\subsection{Local expansion of the covariance}\label{subsec:local-cov-expansion}

This subsection is not on the main line (see the status paragraph at the beginning of the section).

\begin{theorem}[Localization of the fluctuation covariance]\label{thm:local-covariance}
For bonds $b,b'\in\mathfrak F$ let $\mathsf S=\{b_-,b'_-\}$ and $\Lam_m=\Lam_m(\mathsf S)$. Then
\[
C_k(b,b')=\sum_{m\ge0}C_{k,m}(b,b'),\qquad\Delta_k(b,b')=\sum_{m\ge0}\Delta_{k,m}(b,b'),
\]
with $C_{k,0}=C_{\Lam_0}(b,b')$ and $C_{k,m}=C_{\Lam_m}(b,b')-C_{\Lam_{m-1}}(b,b')$ for $m\ge1$, and similarly for $\Delta_k$. Each $C_{k,m}(b,b')$ and $\Delta_{k,m}(b,b')$ depends only on $U$ on $\widehat{\Lam_m^+}$. They are covariant, analytic under complex perturbations, and bounded by
\[
|C_{k,m}(b,b')|+|\Delta_{k,m}(b,b')|\le C'_Le^{-\gamma_2(m+|b-b'|)} .
\]
\end{theorem}

\begin{proof}\hasdetail{thm:local-covariance}
Lemma~\ref{lem:telescoping} with Proposition~\ref{prop:local-fluct}(c) for $Y=\Lam_{m-1}\subset Y'=\Lam_m$. The points $w\in\Sigma_Y$ have distance at least $m-1-R_0-10M-2L$ from $\mathsf S$. So
\[
\sum_{w\in\Sigma_Y}e^{-\gamma''(|b-w|+|w-b'|)}\le C_Le^{-\gamma''(m-R_0')/2}e^{-\gamma''|b-b'|/2},
\]
with $R'_0=R_0+10M+2L+1$. The factor $e^{\gamma''R'_0/2}$ is absorbed into $C'_L$, and $\gamma_2=\gamma''/2$.
\end{proof}

\begin{corollary}[Locality under large changes of the background]\label{cor:large-changes}
Let $U_1$, $U_2$ be two backgrounds satisfying the hypotheses above, which coincide on the bonds of $\hat Y$ for a domain $Y$. Then for bonds $b,b'\in\mathfrak F$ starting in $Y$, with $\delta=\operatorname{dist}(\{b_-,b'_-\},\Tf\setminus Y)$,
\[
|C_k(U_1)(b,b')-C_k(U_2)(b,b')|+|\Delta_k(U_1)(b,b')-\Delta_k(U_2)(b,b')|\le C_Le^{-\gamma_2(\delta-R'_0)} .
\]
\end{corollary}

\begin{proof}\hasdetail{cor:large-changes}
The terms with $\widehat{\Lam_m^+}\subset\hat Y$, i.e.\ $m\le\delta-R'_0$, coincide. The others sum to at most $C_Le^{-\gamma_2(\delta-R'_0)}$.
\end{proof}

Corollary~\ref{cor:large-changes} compares two arbitrary regular backgrounds that agree on a region. They may differ by large amounts elsewhere, which Proposition~\ref{prop:base-dependence} does not cover. For backgrounds of two renormalization group steps, $C_k(U_i)$ is the true fluctuation covariance of Part~VII.

\begin{remark}[Use in the cluster expansion]\label{x:37.9}
Theorem~\ref{thm:local-covariance} is the input for a cluster expansion with an interpolated covariance, e.g.\ the Brydges--Kennedy forest formula. That formula needs only the decay of the kernel and the positivity of the interpolated covariance. The entries of $C_k$ then appear in the activities, and expanding each entry by the theorem makes the dependence on the background local. This localization is of radius type, which is not sufficient for the polymer expansion (Section~\ref{sec:interaction-local}). The main line uses instead the weakened random walks of Part~IX, with the operators of this section as local inverses; Proposition~\ref{prop:local-fluct} is used there at level~(L5). Ba\l aban's random-walk expansion \cite{BalBackground} is another alternative. It can be built from the local inverses $C_Y$, with the parametrix of Section~\ref{sec:random-walk} and the local entries $\Delta_{k,m}$ in the remainder, but it is not needed here.
\end{remark}
